\documentclass[12pt]{article}
\usepackage[utf8]{inputenc}
\usepackage[T1]{fontenc}
\usepackage[margin=1in]{geometry}
\usepackage{setspace}
\usepackage{xcolor}
\usepackage[colorlinks=true, urlcolor=blue!60!black, linkcolor=blue!60!black]{hyperref}
\setlength{\parindent}{0em}
\setlength{\parskip}{0.45em}
\renewcommand{\baselinestretch}{1.0}
\begin{document}
\begin{center}
{\Large \textit{Lemon v.\ Kurtzman}, 403 U.S.\ 602 (1971)\par}
\vspace{0.3em}
{\normalsize Case Brief --- One-Page Edition\par}
\end{center}

\textbf{Memory Jogger.} Parochial-school subsidies fostered excessive entanglement with religion; source of the three-prong \textit{Lemon} test (\href{https://supreme.justia.com/cases/federal/us/403/602/}{Justia}).

\textbf{Facts.} Rhode Island paid salary supplements to teachers of secular subjects in nonpublic, mostly Catholic, schools. Pennsylvania reimbursed such schools for secular teachers' salaries, textbooks, and materials, subject to state audits.

\textbf{Procedural History.} Rhode Island's Act was struck down below; Pennsylvania's challenge was dismissed. The appeals were consolidated.

\textbf{Votes.} 8--0 (Pennsylvania) and 8--1 (Rhode Island): Burger for the Court; Douglas and Brennan concurred; White dissented in part.

\textbf{Holding.} Both statutes violate the Establishment Clause: a law needs a secular purpose, a primary effect neither advancing nor inhibiting religion, and no excessive entanglement.

\textbf{Reasoning.} A teacher in a pervasively religious school, unlike a textbook, cannot be inspected once; enforcing secular-only rules would require continuing surveillance and audits, and annual appropriations would breed religious political division. White called this an ``insoluble paradox.'' \emph{Later:} \textit{Kennedy v.\ Bremerton} (2022) declared \textit{Lemon} abandoned.

\textbf{Examples.} \emph{Same-Side:} policed salary subsidies. \emph{Opposite-Side:} neutral vouchers (\textit{Zelman}). \emph{Fence-Sitter:} public tutors on parochial campuses.

\textbf{Critique.} Originalists call the test ahistorical and manipulable, and Justice Scalia likened it to a horror-movie ghoul that keeps rising from its grave; separationists defend its protection of religious minorities and warn that a history test favors majority faiths.

\textbf{Key Quotation.} ``First, the statute must have a secular legislative purpose; second, its principal or primary effect must be one that neither advances nor inhibits religion, \textit{Board of Education v.\ Allen}, 392 U.S.\ 236, 243 (1968); finally, the statute must not foster `an excessive government entanglement with religion.'\,'' --- Chief Justice Burger, for the Court.

\end{document}
